% GENERATED FILE. Edit canonical lessons, then run npm run generate:books.
% canonical-source-hash: fnv1a64:5e52f3d8f5e5b7ca
% canonical-lessons: IT-C508-allenarsi, IT-C508-allenamento, IT-C508-flessione, IT-C508-addominali, IT-C508-sudare

\chapter{In palestra}
\label{ch:it-in-palestra}
\hlchaptermodality{508}

\begin{chapteropening}
\textbf{By the end of this chapter:} \emph{I can talk about training at the gym.}

Complete the last lesson of chapter 508: I can talk about training at the gym.
\end{chapteropening}

\section[allenarsi]{allenarsi — to train, to work out}

\label{lesson:IT-C508-allenarsi}



\begin{quote}
Before the new one: say the Italian for rake, then the Italian for to hoe.
\end{quote}

\subsection*{You'll want to know: allenarsi}
\textbf{allenarsi} — \textquotedblleft{}to train, to work out\textquotedblright{}.

\begin{quote}
Mi alleno tre volte alla settimana.
\end{quote}

\begin{quote}
\emph{I train three times a week.}
\end{quote}

The coach is \emph{l'allenatore}.

\begin{cousinweb}
\emph{allenarsi} is built on \emph{lena} \textquotedblleft{}vigour, breath\textquotedblright{}, with \emph{a-} and the verb ending; the origin of \emph{lena} is not settled.
\end{cousinweb}

\begin{grammarlens}[title={What you've built}]
One more word for telling, explaining and describing at B1.
\end{grammarlens}

\subsection*{Your turn}
\begin{itemize}
\raggedright
  \item \emph{Say it:} \emph{allenarsi}
  \item \emph{Say it:} \emph{allenarsi}, once more
  \item \emph{Say it:} a sentence of your own with \emph{allenarsi}
  \item \emph{Recall:} say the Italian for rake, then the Italian for to hoe, then say \emph{allenarsi} again
\end{itemize}

\begin{tcolorbox}[breakable,colback=teal!4,colframe=teal!35!black,title={Before you move on}]
What does allenarsi mean? (\textquotedblleft{}To train, to work out\textquotedblright{}.) Say it once more.
\end{tcolorbox}
\section[l'allenamento]{l'allenamento — training, practice}

\label{lesson:IT-C508-allenamento}



\begin{quote}
Before the new one: say the Italian for to hoe, then the Italian for to train.
\end{quote}

\subsection*{You'll want to know: l'allenamento}
\textbf{l'allenamento} — \textquotedblleft{}training, practice\textquotedblright{}.

\begin{quote}
L'allenamento di oggi è stato durissimo.
\end{quote}

\begin{quote}
\emph{Today's training was really hard.}
\end{quote}

\emph{Saltare l'allenamento} — to skip training.

\begin{cousinweb}
\emph{allenamento} is built on \emph{allenare} \textquotedblleft{}to train\textquotedblright{}, from \emph{lena} \textquotedblleft{}vigour, breath\textquotedblright{}, with \emph{-mento} for the action; the origin of \emph{lena} is not settled.
\end{cousinweb}

\begin{grammarlens}[title={What you've built}]
One more word for telling, explaining and describing at B1.
\end{grammarlens}

\subsection*{Your turn}
\begin{itemize}
\raggedright
  \item \emph{Say it:} \emph{l'allenamento}
  \item \emph{Say it:} \emph{l'allenamento}, once more
  \item \emph{Say it:} a sentence of your own with \emph{l'allenamento}
  \item \emph{Recall:} say the Italian for to hoe, then the Italian for to train, then say \emph{l'allenamento} again
\end{itemize}

\begin{tcolorbox}[breakable,colback=teal!4,colframe=teal!35!black,title={Before you move on}]
What does l'allenamento mean? (\textquotedblleft{}Training, practice\textquotedblright{}.) Say it once more.
\end{tcolorbox}
\section[la flessione]{la flessione — push-up; bending}

\label{lesson:IT-C508-flessione}



\begin{quote}
Before the new one: say the Italian for to train, then the Italian for training.
\end{quote}

\subsection*{You'll want to know: la flessione}
\textbf{la flessione} — \textquotedblleft{}push-up; bending\textquotedblright{}.

\begin{quote}
Ogni mattina faccio venti flessioni.
\end{quote}

\begin{quote}
\emph{Every morning I do twenty push-ups.}
\end{quote}

Also \emph{il piegamento}. In grammar, \emph{la flessione} is inflection.

\begin{cousinweb}
\emph{flessione} is a learned borrowing from Latin \emph{flexiō} \textquotedblleft{}a bending\textquotedblright{} (Lewis \& Short), from \emph{flectere} \textquotedblleft{}to bend\textquotedblright{}.
\end{cousinweb}

\begin{grammarlens}[title={What you've built}]
One more word for telling, explaining and describing at B1.
\end{grammarlens}

\subsection*{Your turn}
\begin{itemize}
\raggedright
  \item \emph{Say it:} \emph{la flessione}
  \item \emph{Say it:} \emph{la flessione}, once more
  \item \emph{Say it:} a sentence of your own with \emph{la flessione}
  \item \emph{Recall:} say the Italian for to train, then the Italian for training, then say \emph{la flessione} again
\end{itemize}

\begin{tcolorbox}[breakable,colback=teal!4,colframe=teal!35!black,title={Before you move on}]
What does la flessione mean? (\textquotedblleft{}Push-up; bending\textquotedblright{}.) Say it once more.
\end{tcolorbox}
\section[gli addominali]{gli addominali — abs; sit-ups}

\label{lesson:IT-C508-addominali}



\begin{quote}
Before the new one: say the Italian for training, then the Italian for push-up.
\end{quote}

\subsection*{You'll want to know: gli addominali}
\textbf{gli addominali} — \textquotedblleft{}abs; sit-ups\textquotedblright{}.

\begin{quote}
L'allenatore ci fa fare cento addominali.
\end{quote}

\begin{quote}
\emph{The coach makes us do a hundred sit-ups.}
\end{quote}

The muscles, and also the exercise: \emph{fare gli addominali}.

\begin{cousinweb}
\emph{addominali} is built on \emph{addome} \textquotedblleft{}belly\textquotedblright{}, a learned borrowing from Latin \emph{abdōmen} (Lewis \& Short).
\end{cousinweb}

\begin{grammarlens}[title={What you've built}]
One more word for telling, explaining and describing at B1.
\end{grammarlens}

\subsection*{Your turn}
\begin{itemize}
\raggedright
  \item \emph{Say it:} \emph{gli addominali}
  \item \emph{Say it:} \emph{gli addominali}, once more
  \item \emph{Say it:} a sentence of your own with \emph{gli addominali}
  \item \emph{Recall:} say the Italian for training, then the Italian for push-up, then say \emph{gli addominali} again
\end{itemize}

\begin{tcolorbox}[breakable,colback=teal!4,colframe=teal!35!black,title={Before you move on}]
What does gli addominali mean? (\textquotedblleft{}Abs; sit-ups\textquotedblright{}.) Say it once more.
\end{tcolorbox}
\section[sudare]{sudare — to sweat}

\label{lesson:IT-C508-sudare}



\begin{quote}
Before the new one: say the Italian for push-up, then the Italian for abs.
\end{quote}

\subsection*{You'll want to know: sudare}
\textbf{sudare} — \textquotedblleft{}to sweat\textquotedblright{}.

\begin{quote}
Con questo caldo sudo anche da fermo.
\end{quote}

\begin{quote}
\emph{In this heat I sweat even standing still.}
\end{quote}

The noun is \emph{il sudore}.

\begin{cousinweb}
\emph{sudare} is inherited from Latin \emph{sūdāre} \textquotedblleft{}to sweat\textquotedblright{} (EtymDB; Lewis \& Short).
\end{cousinweb}

\begin{grammarlens}[title={What you've built}]
One more word for telling, explaining and describing at B1.
\end{grammarlens}

\subsection*{Your turn}
\begin{itemize}
\raggedright
  \item \emph{Say it:} \emph{sudare}
  \item \emph{Say it:} \emph{sudare}, once more
  \item \emph{Say it:} a sentence of your own with \emph{sudare}
  \item \emph{Recall:} say the Italian for push-up, then the Italian for abs, then say \emph{sudare} again
\end{itemize}

\begin{tcolorbox}[breakable,colback=teal!4,colframe=teal!35!black,title={Before you move on}]
What does sudare mean? (\textquotedblleft{}To sweat\textquotedblright{}.) Say it once more.
\end{tcolorbox}
